\documentclass[11pt,a4paper]{article}
\usepackage[T1]{fontenc}
\usepackage[utf8]{inputenc}
\usepackage{lmodern}
\usepackage[margin=25mm,headheight=14pt]{geometry}
\usepackage{booktabs,tabularx,array,amsmath,enumitem,fancyhdr}
\usepackage[hidelinks,bookmarksopen=true]{hyperref}
\hypersetup{pdftitle={Rakshak: Final Project Report},pdfauthor={Vishwajeet Singh Bhati; Sumit Kumar; Lakshya Patidar}}
\setlength{\parindent}{0pt}
\setlength{\parskip}{6pt}
\renewcommand{\arraystretch}{1.2}
\setlist{nosep,leftmargin=1.4em}
\setcounter{tocdepth}{1}
\pagestyle{fancy}
\fancyhf{}
\fancyhead[L]{\small Rakshak}
\fancyhead[R]{\small Final Project Report}
\fancyfoot[C]{\thepage}
\renewcommand{\headrulewidth}{0.3pt}
\newcolumntype{Y}{>{\raggedright\arraybackslash}X}
\newcolumntype{L}[1]{>{\raggedright\arraybackslash}p{#1}}
\newcommand{\source}[1]{{\small\textit{Evidence: #1.}\par}}
\newcommand{\chapterpage}[1]{\clearpage\section{#1}}
\newcommand{\repo}{https://github.com/vishwajeet-singh-bhati/rakshak}
\newcommand{\repolink}[2]{\href{https://github.com/vishwajeet-singh-bhati/rakshak/blob/main/#1}{#2}}
\begin{document}

\begin{titlepage}
\centering
\vspace*{16mm}
{\Huge\bfseries RAKSHAK\par}
\vspace{7mm}
{\Large On-device Network Intrusion, SMS Scam\par and Phishing URL Detection\par}
\vspace{12mm}
{\large Final Project Report\par}
{\normalsize Bharat AI-SoC Student Challenge\par}
\vspace{18mm}
{\large\bfseries Team Members\par}
\vspace{4mm}
{\large Vishwajeet Singh Bhati\par Sumit Kumar\par Lakshya Patidar\par}
\vspace{14mm}
{\large\bfseries Faculty Mentor\par}
\vspace{3mm}
{\large Dr. Rohit Kumar Das\par}
\vfill
{\large Indian Institute of Information Technology\par Design and Manufacturing Kurnool\par}
\vspace{8mm}
{\normalsize 25 November 2026\par}
\vspace{3mm}
{\small Application snapshot: 0.18.0-passive-notices, version code 53\par}
\end{titlepage}

\clearpage
\setcounter{page}{2}
\thispagestyle{plain}
\renewcommand{\contentsname}{Index}
\tableofcontents
\vfill
{\small Results are reported against their recorded dataset, device and test date.
Earlier experiments are identified separately from the current application.
This document is prepared for faculty review; it does not record mentor approval or final form submission.}

\chapterpage{Project Overview}
\subsection*{Purpose}
Rakshak is an Android security prototype that brings three checks into one application:
network intrusion detection, SMS scam detection and phishing URL detection.
The aim is to warn users about suspicious activity in language they can understand.
Detection runs locally. A local language model can add a short guidance sentence after the application has already shown a factual warning.

The application observes and warns. It does not block connections, delete messages,
close another application or take recovery action on the user's behalf.
An alert indicates a model estimate or a recorded indicator; it does not establish that a device has been compromised.

\subsection*{Tasks and selected models}
\begin{tabularx}{\textwidth}{@{}L{28mm}YY@{}}
\toprule Task & Model or method & Output \\\midrule
Network monitoring & Binary multilayer perceptron, followed by a CNN/attention/LSTM diagnosis model & Benign or suspicious; estimated attack family when sufficiently confident.\\
SMS screening & Word TF-IDF logistic regression, rules and selective MobileBERT-22 inference & Ham, spam or smishing evidence, combined into an alert decision.\\
Links in messages & Character CNN, character TF-IDF classifier, logistic score combiner and local threat indicators & PHISHING, LIKELY\_SAFE or UNCERTAIN.\\
User explanation & Factual templates plus Llama 3.2 1B Instruct, SpinQuant INT4 & Immediate account of the observation, followed by checked guidance.\\
\bottomrule
\end{tabularx}

\subsection*{Implementation}
The Android application is written in Kotlin. Learned model artifacts run through
ExecuTorch Android 1.5.1, with CPU execution. Python/PyTorch tooling supports training,
preprocessing, export and evaluation. The network path uses rooted packet capture
and a shared feature-extraction implementation. PowerShell and Python scripts operate
the controlled phone tests from a Windows work PC. SQLite stores recent activities.

\subsection*{Main findings}
The system is operational on the test Pixel 6a. Actual phone tests show strong port-scan
responses and successful alerts for the isolated SQL-injection example. They also expose
missed XSS, bot-emulator and resource-denial behaviour, and false alarms during normal
authentication. The SMS model has useful archived results, while benign carrier-message
behaviour still needs work. The URL checker runs quickly, but its available benchmark
does not support an overall accuracy or false-positive claim.

\chapterpage{Application Flow and User Controls}
\subsection*{Three inputs, one explanation path}
\begin{center}
\fbox{\begin{minipage}{0.91\textwidth}
Network packets $\rightarrow$ flow and history features $\rightarrow$ binary gate
$\rightarrow$ optional family diagnosis.\par
Incoming SMS $\rightarrow$ normalize text and extract links $\rightarrow$ SMS cascade.\par
Extracted links $\rightarrow$ independent URL detector.\par
Suspicious result $\rightarrow$ immediate factual note and notification
$\rightarrow$ optional local Llama guidance $\rightarrow$ checked update or default advice.
\end{minipage}}
\end{center}

SMS and URL analysis are separate decisions. A normal-looking message can contain a
risky link. A suspicious message need not contain a link. There is no additional neural
model trained jointly on their outputs; combination rules bring the evidence together.
No URL in a message means no automatic URL-model call for that message.

\subsection*{Controls available to the user}
Each detector has its own switch. Automatic incoming-SMS analysis also requires the
corresponding Android permission and monitoring setting. Users can paste or share text
for a manual check. Activity history is retained for seven days, with a confirmation
before clearing previous records. Alerts show the observation and suggested action.

App attribution uses observed socket ownership where available. A receiving application
is described as the receiver; it is not called the attacker. Closed ports and short-lived
connections may have no identifiable owner. The interface states that limitation rather
than inventing an app or Android service name.

An optional ``expected activity'' action can suppress a narrowly matching future notice
for a limited period, with an undo option. It does not retrain the model or establish
that matching traffic is safe. Monitoring pause requires a user action and confirmation.

\subsection*{Example of the intended experience}
Suppose another lab PC repeatedly tries many ports on the phone. The network model may
identify a possible port scan. The first note explains that many connection attempts were
observed and asks whether the user expected a device search or a connection test.
It names an app only if ownership evidence exists. Llama may rephrase the advice after
the note is visible. A generation failure leaves the factual note in place.

This example describes the flow; it is not an additional test result.
Actual measured results are given in Sections 6--8.
\source{E1, E6, E8, E10}

\chapterpage{Paper Reconstruction and Context Experiment}
\subsection*{Starting point}
The network work began with the paper \textit{A Lightweight Embedded Intelligent Threat
Detection System Using TSANet on Android Platforms} (Malode et al., ISED 2025).
We reconstructed its CNN/LSTM attention design on CIC-IDS2017. The paper did not provide
author code, the precise 32-feature list, saved split indices or sufficient implementation
details for an exact reproduction.

The project therefore used a declared stratified 70/30 training/test split, seed 42,
a fixed 32-feature list, training-only standardization, Conv1D kernel 3 with padding 1,
sigmoid spatial attention and softmax temporal attention. The model had 42,761 parameters.
Adam learning rate 0.001 and batch size 256 were used for 10 epochs, followed by
5 fine-tuning epochs for independent magnitude-pruning branches. Duplicates were retained;
the fixed final checkpoints were evaluated without early stopping.

\subsection*{Connection context}
Isolated flow statistics poorly distinguished some rare attacks. The next experiment
added 78 source/destination history features over 2-, 5- and 15-minute windows.
An equal-capacity control and context model each had 46,841 parameters, the same initial
weights, data order and 10+5 epoch schedule. Only eligible earlier training flows populated
the history; labels and literal addresses were not model inputs.

\begin{center}\small
\begin{tabular}{@{}lrrr@{}}
\toprule Variant & Accuracy (\%) & Macro F1 (\%) & Benign FPR (\%)\\\midrule
Reconstruction & 97.0256 & 76.6327 & 1.68933 \\
Reconstruction, 30\% pruned & 97.3240 & 77.1588 & 2.16379 \\
Matched control & 97.1808 & 76.9501 & 1.53420 \\
Matched control + context & 99.9815 & 98.1405 & 0.00484 \\
Context, 30\% pruned & 99.9850 & 98.6028 & 0.00558 \\
\bottomrule
\end{tabular}
\end{center}

The paper reports 97.35\% baseline test accuracy and 98.30\% after 30\% pruning.
Our corresponding figures were 97.0256\% and 97.3240\%; they do not exactly match.
The missing author details prevent attribution of the difference to one cause.
On the matched unpruned comparison, context increased macro F1 by 21.1904 percentage
points, from 76.9501\% to 98.1405\%.

These are completed-flow, random-row benchmark results from a small set of related
campaigns, with one training seed and retained duplicates. They do not demonstrate
unseen-device or everyday-phone performance. The context method and later phone feature
contract are different experiments; the 99.98\% figure is not the current app's accuracy.
\source{E2}

\chapterpage{Current Network Input and Model Architecture}
\subsection*{What is monitored}
The rooted phone capture supplies IP and transport-layer observations, mainly TCP/UDP
headers, packet lengths, direction and timing. A flow is a group of packets associated
with a connection. Bounded protocol indicators are extracted where the captured data
supports them. The application does not decrypt HTTPS or SSH, inspect another app's
internal state, or turn every network packet into a full application transaction.

The input consists of two FP32 tensors: 32 flow values and 115 context/protocol values,
for a total of 147 features. The flow values describe duration, directional packet and
byte counts, packet-length statistics, packet/byte rates, inter-arrival statistics and
SYN/RST/ACK counts. Context describes recent connections and supported protocol evidence,
with a longest history window of 300 seconds. Literal IP addresses, labels and calendar
times are not fed to the learned classifiers.

Preprocessing is embedded in the exported wrapper. Nonnegative flow values and
18 designated protocol counters receive \texttt{log1p}; the remaining context values
retain their contract-defined form. Training-fitted means and scales standardize the
combined vector. No normalization is fitted on live observations or test rows.

\subsection*{Stage 1: binary gate}
The 18,913-parameter MLP is
$147\rightarrow96\rightarrow48\rightarrow1$, with ReLU hidden activations and a
sigmoid decision score. The deployed gate threshold is 0.5408333333.
A score below it ends the network decision as benign. The gate artifact is 82,820 bytes.

\subsection*{Stage 2: family diagnosis}
Only a positive gate invokes the 173,721-parameter diagnosis model. The 32 flow values
form an ordered feature sequence, not a sequence of 32 packets. Parallel Conv1D-64 and
LSTM-64 branches apply spatial and temporal attention. Batch normalization and fusion
feed an LSTM-128. A separate $115\rightarrow32\rightarrow16$ ReLU branch encodes context.
The final 144 values feed $144\rightarrow128\rightarrow7$ classification layers.
The diagnosis artifact is 949,788 bytes.

The seven outputs are Benign, Bot, Brute Force, DDoS, DoS, Port Scan and Web Attack.
The family cutoff is 0.8. A positive gate without an accepted attack-family diagnosis
is retained as \textit{Unclassified}; even a Benign-favoured diagnosis does not veto
the binary gate. This deliberate policy can retain useful uncertain warnings and
also contributes to false alarms. Llama does not decide these labels.
\source{E1, E3}

\chapterpage{Network Data, Training and Validity Controls}
\subsection*{From published flows to phone-aligned inputs}
The paper experiment retained 2,827,829 CIC-IDS2017 rows: 1,979,480 for training and
848,349 for testing. Cleaning excluded 36 Infiltration and 11 Heartbleed rows and
2,867 rows with nonfinite selected features. Seven family labels grouped related subtypes.

The deployed extractor does not reproduce all completed-flow CSV semantics. The later
phone-aligned experiment therefore processed the five original CIC-IDS2017 packet captures
through the shared feature contract, with PARROT2025 and benign CIC-AAGM2017 captures as
additional normal-traffic sources. Its grouped 70/15/15 split contained 4,956,674 training,
1,188,724 validation and 1,098,929 test views. Views are extracted flow segments, not
independent attack campaigns. The old paper 70/30 split is not reused as this split.

\subsection*{Real-phone augmentation used by the current detector}
Twenty accepted phone sessions were assigned chronologically: 14 training, 3 validation
and 3 test sessions. Requests, responses and isolated effects were recorded independently
of model outputs. This produced 848 unique new training rows, 184 validation rows and
184 test rows. The phone labels cover Benign and Web Attack only; they do not add real
examples for every network family.

New training rows received eight-fold exact replay to avoid being overwhelmed by the
large original corpus. These repeats are declared resampling, not extra independent
examples. Validation and test rows were not resampled. Scaling used the combined
training partition only. The original dataset memberships were retained.

Six training jobs, three seeds per architecture, ran for 40 epochs. Epoch checkpoints
were retained without early stopping. Validation selected \texttt{gate\_s0\_final} and
\texttt{diagnosis\_s1\_best}, together with the thresholds, before final test evaluation.
The selected bundle failed the pilot readiness criterion because benign phone flags
remained too high. It was installed subsequently for explicitly requested experimental
testing; installation did not change that failed criterion into a pass.

\subsection*{Controls and remaining uncertainty}
Capture integrity, packet/connection accounting, independent target effects and
native/offline feature agreement were checked. Unrelated background traffic stayed
unlabelled. Failed or ambiguous recordings were not silently made positive examples.
These controls address several sources of leakage and label error, but do not prove
generalization. The new cohort uses the same phone and lab applications, rare classes
remain small, and earlier test sets have already been examined during development.
\source{E3}

\chapterpage{Network Dataset Results}
\subsection*{Binary detection by the deployed model bundle}
TP and FN count attack-labelled rows; FP and TN count normal-labelled rows.
All rates below are percentages. Precision is $TP/(TP+FP)$, recall is $TP/(TP+FN)$,
and benign FPR is $FP/(FP+TN)$.

\begin{center}\small
\begin{tabular}{@{}lrrrrrrr@{}}
\toprule Cohort & TP & FN & FP & TN & Precision & Recall & FPR\\\midrule
Legacy & 158,724 & 110 & 965 & 939,130 & 99.396 & 99.931 & 0.103 \\
Phone & 62 & 2 & 10 & 110 & 86.111 & 96.875 & 8.333 \\
\bottomrule
\end{tabular}
\end{center}

\subsection*{Family classification on the legacy test views}
These are typed results after the full routing policy. F1 is the harmonic mean of
precision and recall. Typed macro F1 is 94.1080\%.
\begin{center}\small
\begin{tabular}{@{}lrrrr@{}}
\toprule Class & Support & Precision (\%) & Recall (\%) & F1 (\%)\\\midrule
Benign & 940,095 & 99.99 & 99.90 & 99.94 \\
Bot & 640 & 56.45 & 86.09 & 68.19 \\
Brute Force & 2,442 & 99.96 & 99.71 & 99.84 \\
DDoS & 17,840 & 99.18 & 99.29 & 99.24 \\
DoS & 90,084 & 99.93 & 99.99 & 99.96 \\
Port Scan & 47,688 & 99.97 & 99.99 & 99.98 \\
Web Attack & 140 & 89.73 & 93.57 & 91.61 \\
\bottomrule
\end{tabular}
\end{center}

\subsection*{Family classification on the held-out phone sessions}
\begin{center}\small
\begin{tabular}{@{}lrrrr@{}}
\toprule Class & Support & Precision (\%) & Recall (\%) & F1 (\%)\\\midrule
Benign & 120 & 98.21 & 91.67 & 94.83 \\
Web Attack & 64 & 88.41 & 95.31 & 91.73 \\
\bottomrule
\end{tabular}
\end{center}
Phone typed macro F1 is 93.2785\%. Its binary accuracy is 172/184, or 93.48\%,
but 10 of 120 normal rows were flagged. That 8.33\% benign FPR matters more for
everyday use than the aggregate accuracy alone. The frozen model before phone
augmentation detected 0/64 phone attack-labelled rows; augmentation detected 62/64,
while introducing the 10 normal false flags.

The low Bot precision and the lack of phone examples for five attack families limit
what these tables establish. A 99.93\% legacy binary recall cannot be described as
99.93\% detection of attacks on a phone. These are correlated rows from declared
cohorts, not a population estimate of attacks detected or alerts per day.
\source{E3; numerical tables generated directly from the saved evaluation}

\chapterpage{Latest Controlled Tests on the Phone}
\subsection*{Setup and completion}
The 25 November campaign used the Pixel 6a with current app code 53, unchanged detector
weights, the Windows controller and isolated test targets. Captured packets, matching
native flow segments, controller events and target outcomes were saved. Warm-up covered
the 300-second connection history and each completed case included observation time.

Seven of thirteen planned cases completed. The next protocol-control case stopped
because the locked phone did not start its test target. That interruption is inconclusive,
not a classifier miss. The table is frozen to the campaign record updated at
21:38:49 IST on 25 November 2026.

\begin{tabularx}{\textwidth}{@{}L{31mm}L{22mm}Y@{}}
\toprule Case & Flagged / segments & Observed outcome\\\midrule
Benign & 0 / 8 & 8 ordinary HTTP requests; no matched flags. \\
Port scan & 505 / 510 & 503 Port Scan; 2 Unclassified flags. \\
HTTP password guessing & 16 / 32 & 5 Web Attack; 11 Unclassified flags. \\
SQL injection & 2 / 3 & 2 Web Attack flags; 1 control stayed Benign. \\
XSS & 0 / 1 & Script execution confirmed in isolated target; missed. \\
Connection exhaustion & 0 / 16 & 14 resource denials, 2 holds; no matched flag. \\
Bounded HTTP burst & 0 / 128 & 128 requests; no matched flag; harm not established. \\
\bottomrule
\end{tabularx}

\subsection*{What these counts mean}
All seven completed captures passed health and total-packet/native-accounting checks.
The scan made 129 connection attempts, including one open port; retries and segmentation
produced 510 model segments. Its 505 flags therefore do not represent 505 independent
attacks. SQL injection included two independently confirmed effects and one normal control.
The password case included 31 failures and one valid login, so 16/32 is not a pure
attack-recall estimate.

The ordinary HTTP control was quiet for its eight matched connections, but three
outside-test flags occurred in the same window. During the first five cases an experimental
device-local alert setting was enabled. It was disabled before connection exhaustion
and the HTTP burst. Background counts from these scope settings must not be pooled.

Notifications are grouped and may remain visible from earlier activity. A visible
notification count is not the number of matched detections. The target evidence and
matched traffic provide the basis for the results above. The sample sizes do not support
general class recall, an overall TP/FP score, or a claim of 100-variety attack coverage.
\source{E4; the accompanying evidence file preserves the aggregate campaign snapshot}

\chapterpage{Earlier Protocol Tests and Missed Behaviours}
\subsection*{Completed tests from 21 November}
These cases used the same learned network bundle with an earlier application and
explanation policy. They are retained as historical protocol evidence, not new code-53
results. Each listed primary connection produced one native segment in these cases.

\begin{center}\small
\begin{tabular}{@{}lrr@{}}
\toprule Test & Primary connections flagged & Correct attack family\\\midrule
Normal protocol control, first & 3 / 9 & Not applicable\\
FTP password guessing & 29 / 32 & 0 / 32\\
SSH password guessing & 32 / 32 & 6 / 32\\
Restricted bot-behaviour emulator & 0 / 12 & 0 / 12\\
Single-source resource-denial load & 18 / 96 & 0 / 96\\
Three-source resource-denial load & 0 / 96 & 0 / 96\\
Normal protocol control, final & 4 / 9 & Not applicable\\
\bottomrule
\end{tabular}
\end{center}

The resource tests caused denial in a dummy application and then recovery. The
three-source test used distinct source IPs with overlapping activity. It was an
application-capacity test, not a volumetric flood of the institute network.
Availability probes were counted separately: one probe was flagged in each load case.
The probe flag must not be reported as detection of a DDoS load request.

\subsection*{Failure analysis}
Seven of eighteen normal control connections were falsely flagged. All four successful
FTP/SSH control logins were among those flags. Meanwhile, the emulator and distributed
load were missed. This exposes the gap between a high legacy benchmark and real protocol
behaviour: training campaign statistics do not guarantee that application intent is
observable in the available features. SSH content is encrypted. A brief normal login
and a malicious sequence can share similar transport statistics.

The family stage was also weak: FTP positives were Unclassified, and only six SSH
attempts received the Brute Force family. The new HTTP password case receiving Web Attack
or Unclassified does not demonstrate FTP/SSH family recognition. The bot emulator was
restricted test software, not real malware; its failure is still a useful regression
case, but not an estimate of recall against every bot family.

These tests establish reproducible failure cases. The current application should not
be presented as a dependable detector for all seven families merely because their
names appear in its output layer.
\source{E5}

\chapterpage{SMS Detection: Architecture and Data}
\subsection*{Input and routing}
Incoming multipart SMS is assembled before analysis. Normalization extracts links,
email addresses and phone numbers, then creates model text with the prescribed replacements.
Actual URL spans are protected during email masking. Links are checked independently
before deeper SMS classification so that a risky-link warning need not wait for it.

Tier 1 uses word unigram/bigram TF-IDF with 13,545 features and a three-class logistic
regression classifier: Ham (normal), Spam (unsolicited) and Smishing (SMS phishing).
It has 40,638 coefficients/intercepts. The fitted model used balanced classes,
$C=1$ and seed 42. Risk/evidence rules determine whether further inference is needed.

Tier 2 is a 22-layer MobileBERT variant with 22,898,947 parameters, using a hybrid
INT8/FP32 ExecuTorch artifact of 37,250,700 bytes. It receives three int64 tensors
of shape $[1,128]$: token IDs, attention mask and token-type IDs, and returns three
class logits. It began from \texttt{google/mobilebert-uncased}; the final two encoder
blocks were omitted for a documented stability problem. Fine-tuning used four epochs,
learning rate $5\times10^{-5}$, batch size 16, maximum length 128 and seed 42;
validation selected epoch 3.

The final Android policy fuses model and rule evidence. A link-only message skips SMS
semantic escalation but keeps the independent URL decision. A narrow guarded fast path
handles standalone celebration phrases. Model class probabilities are not calibrated
probabilities that a sender is fraudulent.

\subsection*{Recorded dataset split}
\begin{center}\small
\begin{tabular}{@{}lrrrr@{}}
\toprule Split & Ham & Spam & Smishing & Total\\\midrule
Training & 3,309 & 539 & 207 & 4,055\\
Validation & 709 & 113 & 47 & 869\\
Test & 709 & 111 & 49 & 869\\
\bottomrule
\end{tabular}
\end{center}
The main sources were UCI SMS Spam and LSDST 2022, with declared Hinglish seed examples
and a separate NCSU phishing source. Exact text, masked templates, near-duplicates and
declared campaigns formed split groups. A 70/15/15 split followed the external holdout
policy. Another 27,885 external rows were retained, dominated by 27,233 smishing rows;
this is not a balanced estimate of ordinary incoming messages.

Original data files are not redistributed. Archive hashes and group-disjointness records
were verified. Exact retraining still requires those permitted sources and split files.
The original Python policy and the later Android routing fixes are versioned separately.
\source{E6}

\chapterpage{SMS Results and Remaining Errors}
\subsection*{Archived comparison on 869 test messages}
These are supplied, archived neural-classifier evaluations, not a new end-to-end Android
warning benchmark. The deployed batch-one PTQ PTE matches the archived artifact hash.
The test set had already been opened for a previous final comparison.

\begin{center}\small
\begin{tabular}{@{}lrrr@{}}
\toprule Variant & Accuracy (\%) & Smishing P (\%) & Smishing R (\%)\\\midrule
FP32 reference & 97.24 & 77.97 & 93.88 \\
PTQ PTE, batch one (deployed) & 97.12 & 78.57 & 89.80 \\
QAT PTE, batch one & 96.78 & 86.05 & 75.51 \\
\bottomrule
\end{tabular}
\end{center}

\subsection*{Selected PTQ PTE: classwise results}
\begin{center}\small
\begin{tabular}{@{}lrrrr@{}}
\toprule Class & Support & Precision (\%) & Recall (\%) & F1 (\%)\\\midrule
Ham & 709 & 98.88 & 99.58 & 99.23 \\
Spam & 111 & 94.95 & 84.68 & 89.52 \\
Smishing & 49 & 78.57 & 89.80 & 83.81 \\
\bottomrule
\end{tabular}
\end{center}
The selected model's macro F1 was 90.85\%. Its confusion matrix is shown below;
rows are true classes and columns are predicted classes.
\begin{center}\small
\begin{tabular}{@{}lrrr@{}}
\toprule True $\backslash$ predicted & Ham & Spam & Smishing\\\midrule
Ham & 706 & 0 & 3\\
Spam & 8 & 94 & 9\\
Smishing & 0 & 5 & 44\\
\bottomrule
\end{tabular}
\end{center}
Thus 44 of 49 smishing messages were recognized, while 12 other messages were assigned
Smishing. Three of 709 ham messages were falsely classified, giving a 0.423\% ham error
rate in this archive. This is a neural-classifier error rate, not the Android cascade's
real-world warning rate.

\subsection*{Phone checks and unresolved weaknesses}
Integration checks covered four sets of 34 Tier-1 fixtures, 14 neural class/token checks
and 230 authored scenarios. Class agreement on those neural fixtures did not imply
strict numeric parity: the maximum logit difference was 0.449647, exceeding the stated
0.05 criterion. That stricter check remains failed.

Normal carrier and bank-message false alarms were reported during use. The original
messages were lost in the factory reset, and further SMS improvement was put on hold
pending new examples. The archived figures do not resolve those reports. Current evidence
does not establish dependable performance on Indian transactional messages, mixed languages
or new scam campaigns. A representative labelled benign SMS cohort is the next priority.
\source{E6; PTQ results use the deployed PTE entry, not the pre-export PTQ entry}

\chapterpage{Phishing URL Detection: Architecture}
\subsection*{Input and preprocessing}
The detector receives an address string extracted from a message or supplied manually.
It does not open the website. Normalization validates the address, handles international
hostnames, lowercases scheme/host, removes default ports and fragments, and creates a
model string containing host, optional nondefault port, path and query.

The character input is int64 $[1,256]$, with 97 vocabulary IDs. Long inputs preserve
the first 191 characters, a truncation token and the final 64 characters. Short inputs
are padded. A second FP32 $[1,27]$ tensor describes address lengths, host structure,
digits/symbols, path/query structure, entropy and indicators such as IP hosts, punycode,
userinfo and nondefault ports. The first 17 numeric features receive \texttt{log1p};
saved training means/scales standardize all 27 features.

\subsection*{Learned ensemble}
\begin{tabularx}{\textwidth}{@{}L{31mm}Y@{}}
\toprule Component & Architecture\\\midrule
Character CNN & 97-entry embedding of width 32; parallel 128-channel convolutions with kernels 3, 5 and 7; ReLU and masked max-pooling; 384 character values. A $27\rightarrow32$ numeric branch joins them. A $416\rightarrow64\rightarrow1$ head returns a logit. Total: 92,577 parameters.\\
Character linear model & Character n-grams of lengths 2--5; 300,000-term TF-IDF vocabulary, logarithmic term frequency and L2 normalization; logistic-regression logit.\\
Learned combiner & $z=0.331740\,z_{CNN}+0.391087\,z_{ngram}+0.037394$; sigmoid score $p$.\\
\bottomrule
\end{tabularx}

The frozen decision bands are PHISHING for $p>0.6831367399$ and LIKELY\_SAFE for
$p<0.0953727082$; intermediate scores are UNCERTAIN. A valid threat-list match can
override the learned result. Shorteners and selected ambiguous cases receive uncertainty
rather than an assumed safe verdict. HTTPS is not a safety guarantee.

\subsection*{Deployment and provenance}
The Android CNN runs through ExecuTorch. The recorded CNN PTE is 381,444 bytes and the
linear-model binary is 6,391,214 bytes. Available source covers parsing, data preparation,
training, calibration and export, but the handoff does not establish the complete original
training/split provenance. Exact whole-model retraining is therefore not demonstrated.
The public repository ships an empty, stale-marked threat database to avoid redistributing
restricted feed contents; measured phone list-assisted results used a dated local snapshot.
\source{E7}

\chapterpage{URL Benchmark and What It Establishes}
\subsection*{Measured phone run}
On 24 November, all 2,048 unique supplied URLs were processed sequentially on the physical
Pixel 6a with zero input/inference failures. No destination was opened. The full detector
returned 325 PHISHING, 1,581 LIKELY\_SAFE and 142 UNCERTAIN results. Of the 325 phishing
decisions, 271 came from saved threat indicators and 54 from the learned ensemble.

The input file exactly repeated the first 2,048 inputs in the supplied export-reference
file. Those reference outcomes are predictions, not truth labels. Original training overlap
is unknown. This benchmark is useful for runtime and supported-label checks, but it is
not a fresh, independently labelled holdout.

\subsection*{Independent labels and known-positive results}
Exact complete-URL matches in a dated PhishTank feed supported 218 phishing labels.
The other 1,830 URLs remained unknown; there were zero independently confirmed benign
labels. Popular-site identity, appearance and absence from a threat feed were not used
to invent negative labels.

\begin{center}\small
\begin{tabular}{@{}lrrrr@{}}
\toprule View on 218 positives & Phishing & Uncertain & Likely safe & Detection rate\\\midrule
Full local pipeline & 218 & 0 & 0 & 100.00\%\\
Learned ensemble before overrides & 172 & 30 & 16 & 78.90\%\\
\bottomrule
\end{tabular}
\end{center}
All 218 full-pipeline positives were saved-list hits. That 100\% is a known-list result,
not model accuracy. Counting uncertainty as a miss, the learned ensemble detected
172/218 and missed 46/218, giving 78.90\% recall and a 21.10\% false-negative rate
on the supported-positive subset. The 107 other PHISHING outputs have unknown truth.

\subsection*{Metrics that cannot be claimed}
Overall accuracy, precision, F1, specificity and false-positive rate cannot be calculated
for the full file without verified benign labels. The 1,581 likely-safe outputs cannot
all be called true negatives. In addition, an empty public threat database will not
reproduce the phone's 218 saved-list hits without a permitted matching feed snapshot.

The detector-only run took 9.417 seconds. Per-URL mean latency was 4.215 ms, median
3.940 ms and p95 5.307 ms; the first URL took 703.793 ms including cold initialization.
These exclude SMS delivery, notification and Llama. A future balanced, fresh and
independently labelled URL test set is needed to measure both detection and false alarms.
\source{E7}

\chapterpage{Local Explanations and Output Checks}
\subsection*{Why the explanation path changed}
Generating an entire warning with Llama was slow and often failed factual checks.
The adopted design first builds the observation, uncertainty, ownership caveat and
manual advice from recorded facts. Llama generates only a short additional guidance
sentence. The model is Llama 3.2 1B Instruct, SpinQuant INT4 EO8, with
1,135,951,488 bytes of model weights, approximately 1.058 GiB.

The worker runs locally in an isolated process. Verified weights can remain warm for
about 120 seconds of inactivity; generation context resets for each activity.
This reuses weights, not another incident's conversation. The EO8 filename describes
embedding/output precision and is not the custom E8 lattice technique reviewed later.

\subsection*{User-facing structure}
\begin{description}[style=nextline,leftmargin=0pt]
\item[What happened] State the suspected pattern, measured activity and whether an app
was identified as sender or receiver. Avoid asserting a confirmed compromise.
\item[What to do] Suggest a simple manual check appropriate to the evidence, such as
reviewing unfamiliar activity or avoiding a suspicious link.
\item[What remains unknown] Explain missing ownership, incomplete capture/history or an
uncertain diagnosis when these affect interpretation.
\end{description}
Raw messages, OTPs and full URLs are not sent into the guidance prompt. App names and
packet counts are rendered from facts, rather than supplied for the language model to
guess. Unsupported, incomplete or conflicting generated advice is rejected, leaving
the factual note and default guidance available.

\subsection*{Recorded live-worker validation}
With monitoring active, eight declared evidence requests covered seven flagged
categories. Seven generated guidance sentences were accepted. One Web Attack request
introduced an unsupported sign-in purpose and was rejected; correct default guidance
remained. All eight returned useful notes without runtime errors, and the measured
capture queue-drop count was zero in those requests.

One cold request took 15.522 seconds; seven warm requests had median 3.707 seconds
and range 2.944--3.792 seconds. A separately labelled notification-preview check verified
the factual note before Llama started; posting and querying it took 10 ms. This is not
packet-to-alert latency. There has been no blinded readability study or population
hallucination-rate measurement. Not every displayed note is an accepted Llama generation.
\source{E8}

\chapterpage{Hardware Use and Optimization Results}
\subsection*{Measurement environment}
Phone measurements used a Pixel 6a, Google Tensor, eight online CPU cores and Android 17.
Accessible RAM was approximately 5.46 GiB. ExecuTorch Android 1.5.1 used CPU XNNPACK;
no model GPU acceleration or dedicated VRAM result is claimed. Training and tooling
ran on the Windows/WSL work PC, whose GPU inventory included an RTX 2000 Ada, 16 GB.

\begin{tabularx}{\textwidth}{@{}L{55mm}Y@{}}
\toprule Measurement & Recorded result and scope\\\midrule
Current detector graph & Paired timing experiment: median of block medians 1.583 ms, 5,040 calls. Excludes packet capture and flow/history waiting.\\
Hoisted FP32 graph candidate & 1.182 ms under the same paired experiment: 25.3\% lower median latency; zero family/binary changes in checked inputs. Candidate, not the shipped diagnosis artifact.\\
Old full-note Llama worker & Warm median 38.254 s and cold median 54.605 s, three requests each; none of the six drafts passed the old checks.\\
Adopted short-guidance worker & Warm median 3.707 s (seven), cold 15.522 s (one), in a later installed-worker session. Different sessions and workloads; not a randomized paired speedup.\\
Paired component explanation test & Full-note 21.053 s versus fixed facts + short guidance 3.888 s warm median; 81.5\% less time under that component setup.\\
Resident memory & Earlier monitoring-only app median PSS 51.9 MiB (four samples); production 1B Llama worker sampled peak PSS 1.338 GiB. Different scopes, not a summed universal footprint.\\
CPU and temperature & Earlier warm worker phase medians about 303--306\% CPU (roughly three cores); observed battery temperatures 39.4--41.9$^{\circ}$C in that production session.\\
\bottomrule
\end{tabularx}

\subsection*{Options tested but not adopted}
Custom E8 post-training reconstructions changed 278/1,125 binary decisions and
577 or 655 family decisions, depending on reconstruction; they were rejected.
No packed 2-bit Android Llama kernel was produced. A smaller 135M language model was
faster but failed the additional output-quality screen, so it was not promoted.
One-token advance warm-up moved loading earlier but did not establish a total-time
benefit; the on-demand warm worker was retained.

The practical improvement is an immediate useful note, less generated text and reused
verified weights. The 1B model still needs substantial RAM. Sampled CPU/GPU/battery
counters are workload-dependent; they do not establish energy per alert or battery life.
\source{E8, E9}

\chapterpage{Corrections, Privacy and Present Limitations}
\subsection*{Corrections made during phone use}
Normal-looking activity exposed several problems beyond model weights. Link-only
messages were incorrectly escalated by the SMS path; routing was separated. URL spans
containing email-like text were damaged during masking; extraction order was corrected.
Historical reset notices were being reposted when the app reopened; the current release
refreshes only relevant visible notices. Reset wording now explains that Rakshak observed
an unsuccessful connection and did not block it.

Two alerts reported around a Google sign-in were loopback connections, not recorded
external Google connections. An experimental device-local alert setting had remained on.
Normal external-only scope was restored while preserving history and raw model evidence.
The actual loopback service owner and connection to the sign-in remain unproven.
This scope fix does not establish that ordinary external false alarms are solved.

\subsection*{Privacy and operational scope}
The current main application has no INTERNET permission. Inference and stored activity
are local. SMS permission applies to new carrier messages; Rakshak does not automatically
read WhatsApp, RCS, browser history or every installed app's text. Test helpers are
separate packages with their own network permissions. Research captures, messages,
credentials and downloaded threat feeds are excluded from the public repository.

No traffic-blocking path was found in the current source review. The network branch
requires root and continuing capture permission; this is not a stock, unrooted Android
deployment. Capture drops, process restarts and incomplete 300-second history can affect
results. App ownership is best-effort; a remote attacking program or exact Android
service is not reliably identifiable from a short connection.

\subsection*{What remains weak}
\begin{enumerate}
\item Network false positives remain, particularly normal authentication. The gate can
warn even when diagnosis favours Benign. Suppressing all such warnings could also lose
true attacks and has not been accepted as a safe general repair.
\item Real XSS, restricted bot exchanges and resource-denial tests were missed. Encrypted
content and application intent may be unavailable; low-volume attacks can resemble normal use.
\item Rare-class support is small. Same-phone session splits, earlier test inspection and
missing fresh device/day holdouts constrain generalization claims.
\item SMS benign-message handling and strict neural numerical parity remain unresolved.
URL negatives and original split provenance are incomplete.
\item Llama can add unsupported advice; checking and fallback reduce this risk without
eliminating it. Cold loading and 1B memory use remain material costs.
\end{enumerate}
\source{E6--E10}

\chapterpage{Reproducibility and Next Work}
\subsection*{Submitted source and documentation}
The public repository is \url{https://github.com/vishwajeet-singh-bhati/rakshak}. It contains the Android application, shared
feature code, available training/preprocessing/export code, component model provenance,
test helpers and aggregate results. The main reproduction guide is
\repolink{REPRODUCE.md}{REPRODUCE.md}. Dataset downloads, private captures, SSH keys,
raw SMS and restricted feed snapshots are intentionally separate.

\begin{tabularx}{\textwidth}{@{}L{33mm}Y@{}}
\toprule Stage & Required work\\\midrule
Build and install & Clone the repository; follow the pinned Android/JDK setup and artifact checks; build the app and install it on an authorized ARM64 phone. Root is required for the network path. Enable each intended detector and required permissions.\\
Verify execution & Run the supplied native model/feature contracts and component fixtures. Allow a complete history warm-up after a network restart. Fixture agreement is a software check, not a new accuracy benchmark.\\
Repeat lab cases & Follow the controlled/protocol guides with isolated targets, explicit phone/source addresses, fresh captures and independent target events. Additional PCs and SSH setup are needed for the three-source protocol test.\\
Reproduce training & Obtain permitted publisher files, check hashes, preserve the documented splits and fit preprocessing on training only. New phone collection produces a new cohort, not the original private recordings.\\
Compare results & Keep raw model decisions, family correctness, policy alerts, notifications and incomplete cases separate. Do not merge tests from different scope settings or app revisions without labelling them.\\
\bottomrule
\end{tabularx}

Exact old numbers require the same inputs, split memberships, feature/model artifacts
and runtime. Public users can reproduce the procedures and included regression checks;
they cannot recover excluded private recordings or restricted feeds by cloning alone.
The original full URL split is incomplete, and the SMS exact source splits must be
obtained separately. Those are explicit reproducibility limits.

\subsection*{Next work}
The priority is a fresh labelled normal-use cohort spanning legitimate authentication,
carrier/bank SMS and new devices/days. Attack collection should vary rates, protocol
forms and session conditions, retaining independent effect labels. A frozen final test
should follow validation-only calibration of warning thresholds. URL testing needs verified
negatives and time-separated unknown campaigns. Robustness and alerts per day should be
measured before claiming dependable unattended protection.

The present contribution is an integrated, inspectable on-device prototype with explicit
successes and failures. Faculty review, final video publication and the final submission
remain separate checkpoints; they are not asserted complete by this report.

\chapterpage{References and Evidence Index}
The numbered records below identify the sources of the tables and implementation claims.
Repository paths are clickable in the PDF. The accompanying
\repolink{docs/final_submission_report/EVIDENCE.json}{EVIDENCE.json} pins source hashes
and retains the aggregate 25 November campaign snapshot. Raw private traffic is not
part of the report.

\begin{description}[style=nextline,leftmargin=0pt,itemsep=5pt]
\item[E1 --- Current network contract]
\repolink{models/current/MODEL_CARD.md}{Network model card} and
\repolink{android/app/src/main/assets/aligned_v2_gate_manifest.json}{gate manifest}.
Model-card application labels retain their historical dates; report code 53 comes from
the current build configuration, while the detector hashes remain unchanged.
\item[E2 --- Paper and context]
Malode et al., \textit{A Lightweight Embedded Intelligent Threat Detection System Using
TSANet on Android Platforms}, ISED 2025, DOI
\href{https://doi.org/10.1109/ISED67359.2025.11404952}{10.1109/ISED67359.2025.11404952}.
\repolink{docs/PAPER_AND_CONTEXT.md}{Reconstruction assumptions} and
\repolink{results/tsanet_context/COMPARISON.md}{matched context results}.
\item[E3 --- Network data and held-out results]
\repolink{docs/DATASETS.md}{Dataset sources and split counts};
\repolink{experiments/phone_capture_v1/DATASHEET.md}{phone capture datasheet};
\repolink{results/phone_capture_training_v1/evaluation.json}{complete saved evaluation}.
\item[E4 --- Latest phone campaign]
\repolink{docs/final_submission_report/EVIDENCE.json}{Seven-case snapshot and source hash},
25 November 2026, 21:38:49 IST. Seven valid cases; the remaining six have no completed
current-campaign results.
\item[E5 --- Historical protocol tests]
\repolink{results/protocol_phone_tests/RESULTS.md}{Completed 21 November tests} and
\repolink{results/protocol_phone_tests/COMPLETED_REVIEW.json}{accounting review}.
\item[E6 --- SMS]
\repolink{components/sms/README.md}{Architecture and training};
\repolink{components/sms/provenance/ARCHIVED_EVALUATION.json}{archived split, classwise
results and artifact hashes}; \repolink{results/sms_alert_review_20261124/README.md}{routing review}.
\item[E7 --- URL]
\repolink{docs/rakshak/architecture/README.md}{URL/SMS layer architecture};
\repolink{components/phishing_url/README.md}{model and feed provenance};
\repolink{results/url_benchmark_20261124/README.md}{labelled-subset benchmark}.
\item[E8 --- Explanation path]
\repolink{results/hybrid_guidance_v1/README.md}{Installed hybrid-guidance validation};
\repolink{results/hybrid_guidance_v1/summary.json}{runtime summary}.
Llama model: \href{https://huggingface.co/meta-llama/Llama-3.2-1B-Instruct-SpinQuant_INT4_EO8}{Meta Llama 3.2 1B SpinQuant INT4 EO8}.
\item[E9 --- Optimization]
\repolink{results/performance_v1/README.md}{Paired timing, resource and E8 evaluation};
\repolink{results/forgesight_pr_review/README.md}{warm-up and token-limit investigation}.
\item[E10 --- Current notice and scope corrections]
\repolink{results/passive_notice_review_20261125/README.md}{Passive-notice review};
\repolink{results/signin_scope_review_20261125/README.md}{local-scope review}.
\end{description}
\end{document}
